\begin{helper}
Fix one blue support label \(B\), one red label \(J\), the terminal
coordinate set \(U\), and the finite branch/transcript interface for the
red-second part of the reveal order.  For a realized pair \((\beta,\tau)\),
write
\[
R_{\beta,J},\qquad P_\tau,\qquad A_\tau,\qquad Z_\tau
\]
for the decoded red support, forced-present set, forced-absent set, and
selected absent set.  Assume the terminal product law, \(0\le p\le1/2\),
the fixed-\((B,J)\) realized-cell geometry, compatibility soundness and
compatibility uniqueness for complete terminal assignments, and explicit
released residual block data
\[
\mathcal R_{\beta,\tau}\subseteq 2^U
\]
for realized pairs.  The released blocks satisfy:
block assignments are residual assignments on \(U\); membership in
\(\mathcal R_{\beta,\tau}\) makes
\[
(A_{\rm res}\setminus Z_\tau)\cup B\cup R_{\beta,J}
\]
compatible with \((\beta,\tau)\); reconstructed compatibility puts
\(A_{\rm res}\) back in \(\mathcal R_{\beta,\tau}\); blocks for distinct
realized pairs are disjoint; and the residual cylinder
\[
\{A_{\rm res}\subseteq U:
P_\tau\setminus(B\cup R_{\beta,J})\subseteq A_{\rm res},\
A_{\rm res}\cap(A_\tau\setminus Z_\tau)=\emptyset\}
\]
is contained in \(\mathcal R_{\beta,\tau}\).

Then the fixed labels \(B,J\) admit one residual owner
\[
\operatorname{owner}_{B,J}:2^U\to
  \{\hbox{branch/transcript pairs}\}\cup\{\bot\}
\]
whose fiber over every realized pair contains the corresponding residual
cylinder, and a residual assignment cannot be owned by two different
realized pairs.  Moreover there is a finite residual leaf type
\(\Lambda_{B,J}\), a leaf map \(\iota(\beta,\tau)\), and nonnegative leaf
masses \(\nu\), normalized by
\[
\sum_{\lambda\in\Lambda_{B,J}}\nu(\lambda)\le1,
\]
such that the leaf map is injective on realized pairs and every realized
pair satisfies
\[
p^{|P_\tau|-u_R-u_B}(1-p)^{|A_\tau|-|Z_\tau|}
  \le \nu(\iota(\beta,\tau)).
\]
\end{helper}
\begin{proof}
First construct the owner from the supplied released blocks.  For a residual
assignment \(A_{\rm res}\), if it lies in at least one realized block
\(\mathcal R_{\beta,\tau}\), choose one such pair and set
\(\operatorname{owner}_{B,J}(A_{\rm res})=(\beta,\tau)\); otherwise set the
owner to \(\bot\).  The disjointness hypothesis on the released blocks makes
this choice independent of the chosen witness: if the same residual
assignment lay in both \(\mathcal R_{\beta,\tau}\) and
\(\mathcal R_{\beta',\tau'}\), then the block-disjointness assertion would
force \(\beta=\beta'\) and \(\tau=\tau'\).  This is exactly the uniqueness
invariant isolated by the fixed-\((B,J)\) released-block interface
\(\noderef{FixedSetHistoryCellFixedBJResidualReleasedBlockData}\).

For a realized pair \((\beta,\tau)\), the residual-cylinder containment
hypothesis says that every assignment satisfying
\[
P_\tau\setminus(B\cup R_{\beta,J})\subseteq A_{\rm res},\qquad
A_{\rm res}\cap(A_\tau\setminus Z_\tau)=\emptyset
\]
lies in \(\mathcal R_{\beta,\tau}\).  By the preceding paragraph, the owner
of such an assignment is \((\beta,\tau)\).  Thus the owner fiber of
\((\beta,\tau)\) contains the whole displayed residual cylinder.  If the
same owner value were also \((\beta',\tau')\), equality of the two values of
one function gives equality of the ordered pairs, hence
\(\beta=\beta'\) and \(\tau=\tau'\).

It remains to build the normalized leaf package.  Apply
\(\noderef{FixedSetHistoryCellFixedBJResidualPartitionData}\) to the
released blocks.  That node converts the terminal Bernoulli product measure
on \(2^U\) into residual assignment masses \(\mu(A_{\rm res})\), released
assignment sets, and released block masses \(\rho_{\beta,\tau}\).  Its
hypotheses are exactly the block data assumed here: subset of \(2^U\),
soundness of reconstructed compatibility, converse membership,
disjointness, and residual-cylinder containment.  It gives
\[
\sum_{A_{\rm res}\subseteq U}\mu(A_{\rm res})\le1,
\]
disjoint released assignment blocks, and for every realized pair
\[
p^{|P_\tau|-u_R-u_B}(1-p)^{|A_\tau|-|Z_\tau|}
  \le \rho_{\beta,\tau}.
\]

Now apply
\(\noderef{FixedSetHistoryCellFixedBJResidualDecisionTreeSupport}\) to this
partition data.  It takes the disjoint released blocks and their masses,
uses the pair \((\beta,\tau)\) itself as the residual leaf label, assigns
that leaf the mass \(\rho_{\beta,\tau}\) for realized pairs and mass \(0\)
otherwise, and invokes the released-cylinder packing bound to obtain total
leaf mass at most \(1\).  Injectivity of the leaf map is equality of ordered
pairs, and the residual domination inequality is exactly the lower bound on
\(\rho_{\beta,\tau}\) supplied by the partition-data node.  This yields the
claimed residual owner and normalized product-tree leaf package.
\end{proof}
